\section{Introduction}
\label{sec:introduction}

A well-calibrated forecaster states 80\% confidence intervals that contain the true
outcome 80\% of the time. Decades of research on human judgment document systematic
failure to achieve this standard: people's confidence intervals are systematically
too narrow, a phenomenon termed overconfidence in interval estimates \citep{Alpert1982,
Lichtenstein1982, Klayman1999, Soll2004}. The practical consequences are severe.
Overconfident interval estimates underlie miscalibrated risk assessments in medicine
\citep{ChristensenSzalanski1981}, finance \citep{BenDavid2013}, engineering
\citep{Flyvbjerg2002}, and strategic planning \citep{Lovallo2003}.

Large language models (LLMs) are now routinely used as decision-support tools in
precisely these high-stakes domains, yet whether their probabilistic forecasts are
calibrated remains an open question \citep{Kadavath2022, Lin2022, Xiong2024}. Prior
work has focused on factual question-answering confidence or binary classification
probability estimates. Virtually no study has examined LLM calibration on interval
estimates for real-world quantitative forecasting, or asked how miscalibration varies
across prediction horizons and application domains.

We address these questions with three pre-registered studies that progressively
expand the scope of LLM calibration assessment. Study 1 establishes a baseline:
are three frontier LLMs overconfident when predicting next-day equity prices with
stated confidence intervals? Study 2 extends to multiple prediction horizons (1 to
22 days) using 100 large-cap equities and over 18,000 scored predictions to ask:
does miscalibration scale systematically with forecast horizon? Study 3 extends
across six heterogeneous prediction domains (equities, commodities, cryptocurrency,
foreign exchange, weather, NBA game scores) to ask: is miscalibration domain-specific,
and can its direction reverse?

The human overconfidence literature consistently shows that miscalibration is
domain-specific: experts are most overconfident in their own fields of ostensible
competence \citep{Klayman1999, Soll2004}. It also shows that overconfidence in
interval estimates increases with task difficulty and time horizon \citep{Griffin1992}.
If LLMs inherit domain and horizon structure from their training data, they may
replicate or amplify these asymmetries.

Our results confirm and extend the human literature. Across all three studies,
GPT-4 is the most overconfident model and Claude the most calibrated. Study 2
reveals a striking model divergence at longer horizons: Claude and GPT-4 become
increasingly miscalibrated as the prediction window grows, while Gemini maintains
near-target coverage even at 22 days. Study 3 uncovers a domain-confidence
inversion---all models are overconfident in structured financial markets
(equities, commodities) but underconfident in speculative markets (cryptocurrency,
forex)---a pattern absent from prior accounts of LLM calibration.

This paper contributes to three streams of literature. First, we extend the human
overconfidence literature \citep{Soll2004, Griffin1992} to a new class of forecasting
agent across the largest empirical sample to date. Second, we contribute to the LLM
calibration literature \citep{Kadavath2022, Xiong2024} by focusing on interval
estimates rather than point probabilities and by using real-world outcomes rather than
knowledge-base queries. Third, we provide actionable guidance for practitioners
deploying LLMs in financial and operational decision-support roles.
